% =============================================================================
% SimVLA -- Method Section (LaTeX)
% =============================================================================

\section{Method}
\label{sec:method}

\subsection{Overview}

We present \textbf{SimVLA}, a Vision-Language-Action (VLA) model for robotic manipulation
that couples a pretrained vision-language model (VLM) with a flow-matching action
transformer. Given a language instruction and multi-camera observations, SimVLA predicts
a chunk of future actions in an end-to-end fashion. The model consists of three
components: (1) a VLM backbone for visual and language grounding, (2) a diffusion-style
action transformer that decodes actions conditioned on VLM features, and (3) a
flow-matching training objective that enables efficient, high-quality action generation.

% -----------------------------------------------------------------------------
\subsection{Vision-Language Backbone}
\label{sec:backbone}

We adopt \textbf{SmolVLM-500M-Instruct}~\cite{smolvlm} as our vision-language backbone.
The model encodes up to $V{=}3$ camera views jointly alongside the natural-language task
instruction. Each image is resized to $512{\times}512$ pixels via bicubic interpolation
and normalized with ImageNet statistics
($\boldsymbol{\mu}{=}[0.485, 0.456, 0.406]$,
 $\boldsymbol{\sigma}{=}[0.229, 0.224, 0.225]$).
Multiple views are concatenated along the sequence dimension before being fed to the
model, producing a unified visual-language feature sequence
$\mathbf{z} \in \mathbb{R}^{T_{\text{vlm}} \times d_{\text{vlm}}}$,
where $d_{\text{vlm}}{=}576$.

During early training, the VLM backbone parameters are frozen (for the first $10^{3}$
steps) to allow the action head to bootstrap stable gradients before joint fine-tuning
commences.

% -----------------------------------------------------------------------------
\subsection{Action Transformer}
\label{sec:action_transformer}

The action transformer takes noisy actions, robot proprioception, and VLM conditioning
features as input and predicts the denoising velocity field, following the flow-matching
paradigm.

\paragraph{Input encoding.}
At each denoising step, let
$\mathbf{x}_t \in \mathbb{R}^{T_a \times d_a}$
denote the noisy action chunk, where $T_a{=}10$ is the number of predicted action steps
and $d_a{=}7$ for the LIBERO joint-space action
(end-effector $\Delta x, \Delta y, \Delta z, \Delta\text{roll}, \Delta\text{pitch},
\Delta\text{yaw}$, and gripper command).
The proprioceptive state
$\mathbf{s} \in \mathbb{R}^{d_s}$
($d_s{=}8$: end-effector position, orientation, gripper state)
and a sinusoidal time embedding
$\mathbf{e}_t \in \mathbb{R}^{32}$
are tiled along the action-sequence dimension and concatenated with $\mathbf{x}_t$:
%
\begin{equation}
  \mathbf{h}_0
  = \operatorname{Linear}\!\bigl([\mathbf{x}_t \,\|\, \tilde{\mathbf{s}} \,\|\,
    \tilde{\mathbf{e}}_t]\bigr)
  \in \mathbb{R}^{T_a \times d_h},
  \label{eq:input_enc}
\end{equation}
%
where $d_h{=}768$ is the transformer hidden size and $[\cdot\|\cdot]$ denotes
channel-wise concatenation.

\paragraph{VLM conditioning.}
The VLM output $\mathbf{z}$ is projected to $d_h$ dimensions via a linear layer and
appended to the action token sequence, yielding a joint sequence of length
$T_a + T_{\text{vlm}}$:
%
\begin{equation}
  \mathbf{H}
  = \bigl[\mathbf{h}_0 \,\|\, \mathbf{W}_z\mathbf{z}\bigr] + \mathbf{E}_{\text{pos}},
  \label{eq:vlm_cond}
\end{equation}
%
where $\mathbf{E}_{\text{pos}}$ is a learned positional embedding.

\paragraph{Transformer blocks.}
$\mathbf{H}$ is processed by $L{=}12$ standard Pre-LayerNorm transformer blocks, each
with $N_h{=}12$ attention heads (head dimension 64), MLP expansion ratio 4
($d_{\text{mlp}}{=}3072$), GELU activation, and dropout 0.1.
After the final layer, only the $T_a$ action positions are retained and projected by a
linear head to produce the velocity prediction
$\hat{\mathbf{v}} \in \mathbb{R}^{T_a \times d_a}$.

% -----------------------------------------------------------------------------
\subsection{Flow Matching Training Objective}
\label{sec:flow_matching}

We train the action transformer using the \textbf{Conditional Flow Matching}
(CFM)~\cite{lipman2022flow} framework.
Given a ground-truth normalized action chunk
$\mathbf{a} \in \mathbb{R}^{T_a \times d_a}$
and Gaussian noise $\boldsymbol{\epsilon} \sim \mathcal{N}(\mathbf{0}, \mathbf{I})$,
we define the linear interpolant:
%
\begin{equation}
  \mathbf{x}_t
  = t\,\boldsymbol{\epsilon} + (1-t)\,\mathbf{a},
  \qquad
  t \sim \operatorname{Beta}(1.5,\;1.0),
  \label{eq:interpolant}
\end{equation}
%
where $t$ is clipped to $[0.001, 0.999]$.
The regression target (velocity field) is:
%
\begin{equation}
  \mathbf{u}_t = \boldsymbol{\epsilon} - \mathbf{a}.
  \label{eq:velocity_target}
\end{equation}
%
The training loss is the mean squared error between the predicted and target velocity
fields:
%
\begin{equation}
  \mathcal{L}_{\text{flow}}
  = \mathbb{E}_{t,\boldsymbol{\epsilon}}
    \bigl\|
      \hat{\mathbf{v}}(\mathbf{x}_t, t, \mathbf{z}, \mathbf{s}) - \mathbf{u}_t
    \bigr\|_2^2.
  \label{eq:flow_loss}
\end{equation}

\paragraph{Action normalization.}
Prior to training, all action dimensions are normalized to zero mean and unit variance
using dataset statistics:
%
\begin{equation}
  \mathbf{a}_{\text{norm}}
  = \frac{\mathbf{a} - \boldsymbol{\mu}}{\boldsymbol{\sigma} + \varepsilon},
  \qquad \varepsilon = 10^{-6}.
  \label{eq:norm}
\end{equation}
%
Proprioceptive states are normalized identically.

% -----------------------------------------------------------------------------
\subsection{Adaptive Action Chunking}
\label{sec:adaptive_chunking}

To focus learning on critical transition moments (\emph{e.g.}, contacts and direction
reversals), we optionally apply \textbf{Adaptive Action Chunking} by reweighting the
flow-matching loss with per-step change rates.
For a ground-truth action chunk $\mathbf{a} \in \mathbb{R}^{T_a \times d_a}$, we compute
the normalized step-wise change rate:
%
\begin{equation}
  r_\tau
  = \frac{\|\mathbf{a}_{\tau+1} - \mathbf{a}_\tau\|_2}
         {\max_{\tau'}\|\mathbf{a}_{\tau'+1} - \mathbf{a}_{\tau'}\|_2},
  \qquad \tau = 1, \ldots, T_a - 1.
  \label{eq:change_rate}
\end{equation}
%
Each step is assigned a loss weight $w_\tau = 0.5 + r_\tau \in [0.5, 1.5]$, so the
weighted flow-matching loss becomes:
%
\begin{equation}
  \mathcal{L}_{\text{flow}}^{*}
  = \mathbb{E}_{t,\boldsymbol{\epsilon}}
    \bigl\|
      \sqrt{w_\tau}\,
      \bigl(\hat{\mathbf{v}}_\tau - \mathbf{u}_{t,\tau}\bigr)
    \bigr\|_2^2.
  \label{eq:weighted_flow_loss}
\end{equation}

An auxiliary \textbf{chunk boundary head} --- a two-layer MLP
($d_h \!\to\! d_h/2 \!\to\! 1$, SiLU activation, zero-initialized output layer)
applied to the action features --- regresses the stop-gradient change rate $r_\tau$:
%
\begin{equation}
  \mathcal{L}_{\text{boundary}}
  = \operatorname{MSE}\!\bigl(\hat{r}_\tau,\; \operatorname{sg}(r_\tau)\bigr),
  \label{eq:boundary_loss}
\end{equation}
%
where $\operatorname{sg}(\cdot)$ denotes the stop-gradient operator.
The total loss is:
%
\begin{equation}
  \mathcal{L}
  = \mathcal{L}_{\text{flow}}^{*} + \lambda\,\mathcal{L}_{\text{boundary}},
  \qquad \lambda = 0.1.
  \label{eq:total_loss}
\end{equation}
%
Disabling this option recovers the standard uniform-MSE objective
(Eq.~\ref{eq:flow_loss}), remaining fully backward compatible.

% -----------------------------------------------------------------------------
\subsection{Inference}
\label{sec:inference}

At inference time, actions are generated by solving the learned flow ODE with 10 Euler
steps. Starting from $\mathbf{x}_1 \sim \mathcal{N}(\mathbf{0}, \mathbf{I})$:
%
\begin{equation}
  \mathbf{x}_{t-\Delta t}
  = \mathbf{x}_t
    + \Delta t \cdot \hat{\mathbf{v}}(\mathbf{x}_t, t, \mathbf{z}, \mathbf{s}),
  \qquad \Delta t = \tfrac{1}{10}.
  \label{eq:euler}
\end{equation}
%
The resulting $\mathbf{x}_0$ is denormalized using the stored dataset statistics to
recover the final action chunk, which is sent to the robot controller at a replanning
interval of 5 steps.

% -----------------------------------------------------------------------------
\subsection{Training Details}
\label{sec:training}

We train SimVLA with the \textbf{AdamW} optimizer
($\beta_1{=}0.9$, $\beta_2{=}0.95$, weight decay $= 0$, gradient clipping $= 1.0$).
The learning rate follows a linear warmup over 2{,}000 steps to $10^{-4}$, followed by
cosine decay to $10^{-5}$.
We use a batch size of 32 and train for up to $10^{6}$ iterations using
HuggingFace Accelerate~\cite{accelerate} for multi-GPU distributed training.
Checkpoints are saved every 50{,}000 steps.

\begin{table}[h]
  \centering
  \caption{SimVLA architecture and training hyperparameters.}
  \label{tab:hparams}
  \begin{tabular}{lc}
    \toprule
    \textbf{Hyperparameter} & \textbf{Value} \\
    \midrule
    VLM backbone            & SmolVLM-500M-Instruct \\
    Image resolution        & $512 \times 512$ \\
    Number of camera views  & 3 \\
    Action horizon $T_a$    & 10 \\
    Action dimension $d_a$  & 7 \\
    Proprioception dim $d_s$& 8 \\
    Transformer hidden size $d_h$ & 768 \\
    Transformer depth $L$   & 12 \\
    Attention heads $N_h$   & 12 \\
    MLP expansion ratio     & 4 \\
    Time embedding dim      & 32 \\
    Flow time distribution  & Beta(1.5, 1.0) \\
    Inference steps         & 10 (Euler) \\
    Replanning interval     & 5 steps \\
    \midrule
    Optimizer               & AdamW \\
    Learning rate           & $10^{-4}$ \\
    LR schedule             & Linear warmup + cosine decay \\
    Warmup steps            & 2{,}000 \\
    Batch size              & 32 \\
    Gradient clipping       & 1.0 \\
    Training iterations     & $10^{6}$ \\
    Boundary loss weight $\lambda$ & 0.1 \\
    \bottomrule
  \end{tabular}
\end{table}
